\section{Conclusion}\label{sec:conclusion}
\sys\ shows that a 1.5B-parameter open model, adapted with LoRA on a single 6\,GB consumer GPU in about an
hour, reaches 0.733 accuracy on held-out patients and cites the same evidence sentences as expert annotators
far more precisely than its untuned starting point. Every component runs locally, so patient notes never leave
the machine.

Each inference mechanism has its own lesson:
\begin{itemize}
\item \emph{KV cache}: removes the quadratic cost of recomputation, and its benefit grows with context length;
\item \emph{online softmax}: removes the quadratic memory of attention, but needs kernel fusion before it also
saves time;
\item \emph{speculative decoding}: is exact, but at batch~1 on launch-bound hardware it only pays off when the
draft is much cheaper per step than the target, not merely smaller.
\end{itemize}
In reduced precision, agreement with the reference is ``identical up to ties within one rounding step'', and a
careful comparison has to say so explicitly.

The weakest areas are the rare labels and reasoning from the absence of a fact. Natural next steps are more
annotated data for \emph{not included} and \emph{not applicable}, label-balanced or contrastive training,
clinician review of the generated explanations, and a larger validation study: our three-seed runs show that
validation macro-F1 varies by about 0.16 between seeds, so single-run differences in the rare classes cannot be
trusted. Oversampling the rare labels, a fine-tuned draft model and a lower LoRA rank were all tested on the
validation split (Section~\ref{sec:res-extra}); none is shown to help beyond noise except that a fine-tuned
draft raises speculative acceptance from 0.59 to 0.85 (still slower than plain decoding).
